%! Author = chtompki
%! Date = 1/14/26
\documentclass[12pt]{amsart}
% Default margins are too wide all the way around. I reset them here
\setlength{\hoffset}{-.75in}
\setlength{\textwidth}{6.5in}
\setlength{\voffset}{-.5in}
\setlength{\textheight}{9.0in}
\usepackage{amsmath}
\usepackage{leftindex}
\usepackage{amssymb}
\usepackage{amsthm}
\usepackage{enumitem}
\usepackage{setspace}

\newenvironment{solution}
{\begin{proof}[Solution]}
{\end{proof}}

\title{Combinatorics\\Assignment 1}
\author{Rob Tompkins}
\date{\today}

\begin{document}
    \maketitle
    \date{\today}
    \begin{onehalfspace}
        % Problem 1
        \noindent\textbf{1.} There are 100 students at a school and three dormitories $A$, $B$, and $C$
        with capacities 25, 35 and 40 respectively
        \begin{enumerate}[label=(\alph*)]
            \item How many ways are there to fill the dormitories?
            \item Suppose that of the 100 students 50 are men and 50 are women and that $A$
                  is an all-men's dorm, $B$ is an all-women's dorm, and $C$ is co-ed. How
                  many ways are there to fill the dormitories.
        \end{enumerate}
        \begin{solution}
            (a) We begin by asking how many different ways are there to fill a a dorm of size 25 from 100 people:
            \begin{displaymath}
                \binom{100}{25} = \frac{100!}{25!(100 - 25)!}
            \end{displaymath}

            Now if we want to fill the three dorms, we start by filling $A$ with the $\binom{100}{25}$ possibilities
            leaving us 75 students to work with for the next dorm, $B$, which we fill with $\binom{75}{35}$
            possibilities. Lastly, we are left with $C$ which has $\binom{40}{40}$ possibilities. So, the number
            of ways to fill the dorms follows as
            \begin{eqnarray*}
                \binom{100}{25}\binom{75}{35}\binom{40}{40} & = & \frac{100!}{25!(100 - 25)!} \cdot \frac{75!}{35!(75 - 35)!} \\
                                                            & = & \frac{100!}{20!25!35!} \\
                                                            & = & 7.0 \times 10^{44} \text{ ways.}
            \end{eqnarray*}

            (b) Let us be under the assumptions stipulated in (b) above with $A$ being all male, $B$ being all female,
            and $C$ being co-ed. Notice, when $A$ is filled the number of remaining men to fill dorm $C$ will be
            $50 - 25 = 25$. Similarly, the number of women left to fill dorm $C$ will be $50 - 35 = 15$. So, when we
            put this all together using the binomial coefficient as we did before, we see that the number of ways to fill
            the dorms follows as:
            \begin{eqnarray*}
                \binom{50}{25}\binom{50}{35}\binom{40}{40} & = & \frac{50!}{25!25!}\cdot\frac{50!}{35!15!} \\
                                                           & = & \frac{50!^2}{25!^2 35!15!} \\
                                                           & = & 2.845 \times 10^{26} \text{ ways.}
            \end{eqnarray*}
        \end{solution}

        % Problem 2
        \textbf{2.} A committee of five people is to be chosen from a club that boasts a
        membership of 10 men and 12 women. How many ways can the committee be formed if
        it is to contain at least two women? How many ways if, in addition, one
        particular man and one particular woman who are members of the club refuse to
        serve together on the committee?
        \begin{solution}
            For the first portion we proceed in cases using both the addition and multiplication rules:
            \begin{itemize}
                \item Case: 0-men, 5-women: $\binom{12}{5} = 792$
                \item Case: 1-man, 4-women: $\binom{10}{1}\binom{12}{4} = 4950$
                \item Case: 2-men, 3-women: $\binom{10}{2}\binom{12}{3} = 9900$
                \item Case: 3-men, 2-women: $\binom{10}{3}\binom{12}{2} = 7920$
            \end{itemize}
            So we add these values up to get the total number of possibilities for filling our committee, which follows
            as
            \begin{displaymath}
                792 + 4950 + 9900 + 7920 = 23,562 \text{ combinations.}
            \end{displaymath}

            For the next portion of the problem we will comute the complement and subtract it from the whole given just
            above. So assume that our pair of people is infact on the committee. We now have three slots to work with,
            where one of the slots must be a women, and we are working with 9-men and 11-women. So our cases for filling
            these slots follow as:
            \begin{itemize}
                \item Case: 0-men, 3-women: $\binom{11}{3} = 165$
                \item Case: 1-man, 2-women: $\binom{9}{1}\binom{11}{2} = 495$
                \item Case: 2-men, 1-woman: $\binom{9}{2}\binom{11}{1} = 396$.
            \end{itemize}
            So we add these, and we see that the number of misconfigured committees follows as:
            \begin{displaymath}
                165 + 495 + + 396 = 1056,
            \end{displaymath}
            and the number of possible ways to fill the committee follows as
            \begin{displaymath}
                23,562 - 1056 = 22,506.
            \end{displaymath}

        \end{solution}

        % Problem 3
        \textbf{3.} A secretary works in a building located nine blocks east and
        eight blocks north of his home. Every day he walks 17 blocks to work.
        \begin{enumerate}[label=(\alph*)]
            \item How many different route are possible for him?
            \item How many different routes are possible if the one block in the easterly,
                  which begins four blocks east and three blocks north of his home is water
                  (and he can't swim)? (Hint: Count the routes that use the block under water).
        \end{enumerate}
        \begin{solution}
            Consider the $U\times V$ taxicab lattice. We let the point $(m,n)$ represent the destination, and $(0,0)$
            represent the starting point of our secretary. We must move $m$ units to the right, and $n$ units up no
            matter what path we take. So the totalk length of the shortest path is $m+n$. Any path is a sequence of $m+n$
            length over the letterss $R$ ``right'' and $U$ ``up'' (e.g. $RUURURRURRU\ldots$). Now notice that the total
            number of paths is the total number of ways to choose which $m+n$ slots in the sequence will be the horizontal
            moves $R$. So we have:
            \begin{displaymath}
                \text{Number of paths } = \binom{m + n}{m} = \frac{(m+n)!}{m!n!}.
            \end{displaymath}

            (a) We now simply set $U = 9$ and $V = 8$ so the number of paths becomes:
            \begin{displaymath}
                \text{Number of paths } = \binom{17}{8} = \frac{17!}{8!9!} = 24,310
            \end{displaymath}
        \end{solution}

        % Problem 4
        \textbf{4.} There are 20 identical sticks lined up in a row occupying 20 distinct places
        as follows: $\vert\vert\vert\vert\vert\vert\vert\vert\vert\vert\vert\vert\vert\vert\vert\vert\vert\vert\vert\vert$ - six of them are to be chosen.
        \begin{enumerate}[label=(\alph*)]
            \item How many choices are there?
            \item How many choices are there if no two of the chosen sticks can be consecutive?
            \item How many choices are there if there must be at least two sticks between each
                  pair of chosen sticks?
        \end{enumerate}
        \begin{solution}
            (a) Clearly for choosing 6 sticks out of a row of 20, the solution follows as the binomial coefficient:
            \begin{displaymath}
                \binom{20}{6} = 38,760
            \end{displaymath}
            (b) For the next portion of the problem we will reduce it to the classical ``stars and bars'' problem. So,
            we will let our unchosen sticks be given as $*$, and our chosen sticks to be $\vert$, so a representation
            of our solution might look like:
            \begin{displaymath}
                * \vert * \vert * * \vert * * * * \vert * * \vert * \vert * *
            \end{displaymath}
            This is essentially a specific case of the problem of distributing $r$ identical balls into $n$ distinct
            boxes. Let us briefly consider the general case here. But let us consider the general case where there must
            be a ball in every box. First put one ball in each box, and then put a false bottom in the boxes to conceal
            the ball in each box. Now it remains to count the ways to distribute without restriction the remaining
            $r - n$ balls into the $n$ boxes. So the answer is
            \begin{displaymath}
                \binom{(r - n) + n - 1}{r - n} = \frac{\left[(r - n) + n - 1\right]!}{(r - n)!(n - 1)!} = \binom{r-1}{n-1}
            \end{displaymath}
            So our answer falls out that we have 7 - boxes and 14 balls (or stars). And the solution follows as:
            \begin{displaymath}
                \binom{14 - 1}{7 - 1} = \binom{13}{6} = 1716
            \end{displaymath}

            (c) For the third portion of the problem we will use the fact that the nunber of solutions to the equation:
            \begin{displaymath}
                x_1 + x_2 + \cdots + x_k = n
            \end{displaymath}
            such that for all $1\le i \le k$, $x_i > 0$. We know that the number of solutions to this is $\binom{n+k-1}{k-1}$. It
            is also clear that this can be reduced to the number $x_i$ balls in $k$ boxes out of $n$ balls. From here
            we set our linear equation with seven variables for our seven bins above, and set all of our variables
            greater than or equal to two. So we have
            \begin{displaymath}
                x_1 + x_2 + \cdots + x_7  = 14
            \end{displaymath}
            with $x_i \ge 2$. Now set $x_i' \ge 0 $ with $x_i' = x_i - 2$. Then our equation becomes:
            \begin{eqnarray*}
                (x_1' + 2) + (x_2' + 2) + \cdots + (x_7' + 2)  & = & 14 \\
                x_1' + x_2' + \cdots + x_7'  & = & 14 - 14 = 0
            \end{eqnarray*}
            So the number of solutions follows as:
            \begin{displaymath}
                \binom{0 + 7 - 1}{7-1} = \binom{6}{6} = 1
            \end{displaymath}
        \end{solution}

        % Problem 5
        \textbf{5.} How many integral solutions of $x_1 + x_2 + x_3 + x_4 = 30$ satisfy
        $x_1\ge 2$, $x_2\ge 0$, $x_3\ge -5$, and $x_4\ge 8$?
        \begin{solution}
            Similarly to (c) in problem 4, we will simply manipulate the inequalities and use the number of solutions
            to the linear combination of non-negative integers computation. So for our problem let $x_1' \ge 0$ with
            $x_1' = x_1 - 2$, $x_3' \ge 0$ with $x_3' = x_3 + 5$, and $x_4' \ge 0$ with $x_4'=x_4 - 8$. Then our equation
            becomes:
            \begin{eqnarray*}
                (x_1'+ 2) + x_2 + (x_3' - 5) + (x_4' + 8) & = & 10 \\
                x_1' + x_2 + x_3' + x_4' & = & 5
            \end{eqnarray*}
            So our calulation for the number of possible solutions follows as:
            \begin{displaymath}
                \binom{25+ 4 - 1}{4 - 1} = \binom{29}{4} = 23,751.
            \end{displaymath}
        \end{solution}

        % Problem 6
        \textbf{6.} A bagel store sells six different kinds of bagels. Suppose you choose 15
        bagels at random. What is the probability that your choice contains at least one bagel
        of each kind? If one of the kinds of bagels is Sesame, what is the probability that
        your choice contains at least three Sesame bagels?
        \begin{solution}
            Again we are going to use the number of solutions for a linear equation. First though, because we want to
            calculate probabilities. Let's get the total possible number of bagel orders of 15 bagles grabbing from 6 types
            randomly. Let ${A, B, C, D, E, S}$ be our types of bagels. So for $A, B, C, D, E, S \ge 0$ and
            \begin{displaymath}
                A + B + C + D + E + S = 15
            \end{displaymath}
            We quickly see that the total number of possible orders follows as:
            \begin{displaymath}
                \binom{15 + 6 - 1}{6 - 1} = \binom{20}{5} = 15,504.
            \end{displaymath}
            Now let's compute the same thing but instead for $A, B, C, D, E, S > 0$. Given the same linear equation
            above we know from earlier that the total number of orders given these constraints are:
            \begin{displaymath}
                \binom{15 - 1}{6 - 1} = \binom{14}{5} = 2002
            \end{displaymath}
            So,
            \begin{displaymath}
                \text{the probability of getting one of each bagel type in the order } = \frac{2002}{15,504}\cdot 100 = 12.913\%
            \end{displaymath}
            We let $S$ bet the sesame type of bagels in our linear combination. So for us to have at least three sesame
            bagels our variable constraints follow as $A, B, C, D, E \ge 0$ and $S \ge 3$. So we let $S' \ge 0$ be
            defined as $S' = S - 3$. Then our linear combination follows as:
            \begin{eqnarray*}
                A + B + C + D + E + (S' + 3) & = & 15 \\
                A + B + C + D + E + S' & = & 12
            \end{eqnarray*}
            So we fall back to our formula that relies on the binomial coefficient:
            \begin{displaymath}
                \binom{12 + 6 - 1}{6 - 1} = \binom{17}{5} = 6188
            \end{displaymath}
            So our probability of having atleast three sesame begals follows as:
            \begin{displaymath}
                \frac{6188}{15,504}\cdot 100 = 39.912\%.
            \end{displaymath}
        \end{solution}

        % Problem 7
        \textbf{7.} Four (standard) dice (cubes with 1, 2,3, 4, 5, 6, respectively, dots
        on their six faces), each of a different color, are tossed, each landing with one
        of its faces up, thereby showing a number of dots. Determine the following probabilities:
        \begin{enumerate}[label=(\alph*)]
            \item The probability that the total number of dots shown is 6
            \item The probability that at most two of the dice show exactly one dot
            \item The probability that each die shows at least two dots
            \item The probability that the four numbers of dots shown are all different
            \item The probability that there are exactly two different numbers of dots shown
        \end{enumerate}
        \begin{solution}
            (a) We can use our number of solutions for linear equations for this problem, but instead
            let $a_1, a_2, a_3, a_4$ represent the die such that $1 \le x_i \le 6$. We need $x_1 + x+2 + x_3 + x_4 = 6$.
            But because each $a_i \ge 1$, set $b_i = a_i - 1$ so that $b_i \ge 0$. Then
            \begin{eqnarray*}
                (b_1 + 1) + (b_2 + 1) + (b_3 + 1) + (b_4 + 1) & = & 6 \\
                b_1 + b_2 + b_3 + b_4 & = & 2
            \end{eqnarray*}
            So the number of non-negative solutions to the equations follows as
            \begin{displaymath}
                \binom{2 + 4 - 1}{4 - 1} = \binom{5}{3} = 10
            \end{displaymath}
            So the probability of landing this dice roll is
            \begin{displaymath}
                P = \frac{10}{1296} = \frac{5}{648}.
            \end{displaymath}

            (b) Let $k$ be the number of dice showing 1. We want $k = 0, 1, 2$
            We begin by chosing which $k$ dice are 1, $\binom{4}{k}$. Each of the other
            $4 - k$ dice can be $2, 3, 4, 5, 6$ which is $5^{4-k}$ ways. Wo we count
            \begin{eqnarray*}
                \sum_{k=0}^3 \binom{4}{k}5^{4-k} & = & \binom{4}{0}5^4 + \binom{4}{1}{5^3} + \binom{4}{2}5^2 \\
                             & = & 625 + 500 + 150 = 1275.
            \end{eqnarray*}
            So our probability at most two dice show 1 is
            \begin{displaymath}
                P = \frac{1275}{1296} = \frac{425}{432}
            \end{displaymath}

            (c) Each die can be $2, 3, 4, 5, 6$ which is five choices per die that are independent. So the total
            number of outcomes is:
            \begin{displaymath}
                5^4 = 625
            \end{displaymath}
            So the probability each die is at least a two is:
            \begin{displaymath}
                P = \frac{625}{1296}
            \end{displaymath}

            (d) Choose 4 distinct faces for 6 with order since the die are coloured. So we
            have $6 \cdot 5 \cdot 4 \cdot 3 = \leftindex_6 P_4 = 360$. So the probability of all
            four nunbers being different is:
            \begin{displaymath}
                P = \frac{360}{1296} = \frac{5}{18}
            \end{displaymath}

            (e) We have two cases that we have to deal with here where the dice are split evenly 2,2
            and where we have a stand alone die and three repetitions call it 3,1.
            \textbf{Case 1.} 3,1 - Choose the two numbers that we're going to get on the die $\binom{6}{2} = 15$.
            Choose which number appears 3 times: 2 ways to do this. Choose which die is the singleton $\binom{4}{1} = 4$.
            So our count follows as $15 \cdot 2 \cdot 4 = 120$.
            \textbf{Case 2.} 2,2 - Choose the two numbers that we're going to get on the die $\binom{6}{2} = 15$. Choose
            which two dice show the first number and the rest show the other number $\binom{4}{2} = 6$. Which gives the
            count $15 \cdot 6 = 90$.

            So now we combine the results of Case 1 and Case 2 with the total ways to get just two numbers
            as:
            \begin{displaymath}
                120 + 90 = 210
            \end{displaymath}
            So the probability of getting the same two numbers across the die follows as:
            \begin{displaymath}
                P = \frac{210}[1296] = \frac{35}{216}.
            \end{displaymath}
        \end{solution}
    \end{onehalfspace}
\end{document}
